\par\clearpage
\normalmarginpar
\scanpage{5}
\mspage{\handwritten{\tikz[baseline=(f.base)]{\node[inner sep=1pt] (f) {11bis};\draw (f.north west) .. controls ([yshift=-5pt]f.south west) and ([yshift=-5pt]f.south) .. ([yshift=-5pt]f.south east);}}}
\noindent\handwritten{\tikz[overlay]{\draw (-1.3,.30)--(-1.24,.41);\draw (-1.14,.41)--(-.99,.29);}}
\handwritten{\tikz[overlay]{\draw (.72\linewidth,.35) .. controls (.73\linewidth,.23) and (.732\linewidth,.23) .. (.733\linewidth,.35);}}
% The three large cancellation strokes cross the two numbered sections.
\figure{\begin{tikzpicture}
\node[anchor=north west,inner sep=0pt] (body) at (0,0) {\begin{minipage}{.93\linewidth}
\small
\noindent
\deleted{1.3.{\deleted{\unsure{2|3}}}{\unsure{3|2}}.}\quad
\deleted{\inserted{La terminologie de 1.3.1 {\deleted{et 1.3.2}} s'étend aussitôt au cas où}}
\deleted{\deleted{Revenons au cas {\unsure{où|ou}} on se donne deux Groupes $P$ et $Q$ sur $G$}}
\deleted{on a un {\uline{bitorseur}} {\deleted{\unsure{osus|sous}}} {\inserted{$E$}} sous{\deleted{$(G$}} $(G_{P},G_{P})$, {\deleted{resp{\unsure{.|,}} sous $({G\unsure{,}}_{P\mathrm{x}Q},G_{P\mathrm{x}Q})$,}} étant entendu qu'une rigidification du bitorseur $E$ doit être un isomorphisme de $e_{P}(E)$ avec le b{\deleted{o}}it{\deleted{r}}o{\deleted{o}}rs{\deleted{\unsure{e|u}}}eur trivial $\,{}_{s}G_{d}$, {\deleted{donc}} {\inserted{i.e.}}{\handwritten{\tikz[baseline]{\draw (0,-.04) .. controls (.035,.14) and (.16,.06) .. (.19,.25);}}} est défini par une section $e_{E}$ de $e_{P}(E)$ qui est {\uline{centrale}}, i.e. qui satisfait à la condition $g.e_{E}=e_{E}.g\quad(g\mathrel{\handwritten{\unsure{\in|\epsilon}}}G(S)$ . On étend de même la terminologie des birigidifications au cas des bitorseurs.}

\noindent\deleted{1.3.4.}\quad
\deleted{Revenons au cas où on se donne deux Groupes $P$ et {\deleted{$Q$}} {\handwritten{\inserted{$G$}}}, et choisissons la section unité $e_{P}$ pour définir des rigidifications de bitorseurs sur $P$, et des birigidifications de bitorseurs sur $P\mathrm{x}P$. On a des foncteurs naturels}

\begingroup
\csname tagsleft@true\endcsname
\diagram{\[
\begin{tikzcd}[ampersand replacement=\&,column sep=1.7em]
\deleted{\underline{\mathrm{Ext}}(P;\unsure{G'|G})} \arrow[r,"{\handwritten{i}}"]
\& \deleted{\underline{\mathrm{Bitorsrig}}(P,G)} \arrow[r,"{\handwritten{\rlap{\raisebox{.35ex}{\hspace{-.12em}\(\color{gray!55}\scriptstyle\delta\)}}\delta}}"]
\& \deleted{\underline{\mathrm{Bitorsbirig}}(P,P;G){,}}
\end{tikzcd}
\tag*{\deleted{(1.3.4.1)}}
\]}
\endgroup

\noindent\deleted{où les trois expressions écrites désignent respectivement la catégorie des extensions de {\handwritten{\tikz[baseline,overlay]{\node[anchor=south west,inner sep=0pt] at (-.15,.31) {\inserted{Groupes de}};\draw (-.08,-.12) .. controls (-.08,.24) and (.06,.26) .. (.35,.26) .. controls (.75,.26) and (1.15,.24) .. (1.35,.32) .. controls (1.40,.35) and (1.40,.43) .. (1.32,.42);}}}$P$ par $G$, la catégorie des {\deleted{b}}{\handwritten{B}}itorseurs sous $(G_{P},G_{P})$ rigidifiés pour{\textasciigrave} $e_{P}$, et la catégorie des {\deleted{t}}b{\deleted{o}}it{\deleted{r}}orseurs sous$(G_{P\mathrm{x}P},G_{P\mathrm{x}P})$ birigidifiés pour $(e_{P},e_{P})$. On {\deleted{prend}} définit le foncteur $i$ en associant à l'extension $E$ de $P$ par $G$ le {\deleted{torseu}} bitorseur correspondant sur $P$ (1.1.6), rigidifié par la sectio{ń} unité de $E$. O{\deleted{\unsure{n}}}{\handwritten{n}} {\unsure{definit|définit}} le foncteur $\handwritten{\delta}$ par la formu{\deleted{o}}le}
\end{minipage}};
\begin{scope}[x={(body.north east)},y={(body.south west)}]
\draw (.44,.11) .. controls (.35,.27) and (.00,.66) .. (-.047,.725)
 .. controls (-.053,.736) and (-.054,.746) .. (-.043,.75);
\draw (.82,.078) .. controls (.73,.30) and (.50,.69) .. (.31,1.035);
\draw (.91,.505) .. controls (.915,.62) and (.78,.86) .. (.655,1.03);
\end{scope}
\end{tikzpicture}}
% archive sheet number: 4
